\section{RF/Microwave Measurements}
{\color{sub}\footnotesize 6 hrs \gap$\cdot$\gap asked in \mk{22 of 24 papers}
\gap$\cdot$\gap 9 syllabus sub-topics \gap$\cdot$\gap \mk{181 marks, 9.5\% of the archive}
$\rightarrow$ 149 of those 181 marks are \textbf{power measurement}, and the tool is picked
by one number: how many milliwatts.}

\vspace{3pt}\noindent
\colorbox{gdL}{\makebox[\dimexpr\textwidth-2\fboxsep][l]{%
  \textcolor{gd}{\bfseries\footnotesize READ THIS FIRST}}}
\par\nobreak\vspace{3pt}

Chapter 8 is descriptive and it repeats harder than any other chapter: five short-note
titles (bolometry, double-channel bolometer, circulating calorimetry, spectrum analyzer,
VNA) and three long-answer shapes cover everything asked in 24 papers.
\begin{enumerate}[label=\arabic*., leftmargin=1.4em, itemsep=1pt]
  \item \textbf{``Choose a proper power measuring tool and explain its working principle''}
    (72 Ma 7.5 mW, 74 Ma airport radar, 73 Bh, 82 Bh BTS) $\rightarrow$ one table,
    $\S$8.2. Say which band the number falls in, name the sensor, then give that sensor's
    working principle from $\S$8.3 or $\S$8.5. \mk{The table is the mark.}
  \item \textbf{Calorimetry} (9 papers) $\rightarrow$ $\S$8.5. Static vs circulating, one
    formula each, and they are the same formula with $m/t$ replaced by $v\,d$.
  \item \textbf{Bolometer bridge} (9 papers) $\rightarrow$ $\S$8.3. Barretter vs thermistor,
    single bridge, its two faults, double bridge that cures them.
\end{enumerate}
Three things worth knowing before the paper starts:
\begin{itemize}
  \item \textbf{Every microwave power meter is a substitution instrument.} It never measures
    RF. It measures how much \textbf{dc or low-frequency} power had to be taken away to keep
    a bridge balanced, or how much a \textbf{temperature} rose. Open any power answer with
    that sentence.
  \item \textbf{The three bands are the whole decision}: low $<$ 10 mW, medium 10 mW to
    10 W, high $>$ 10 W. A radar is quoted in \textbf{peak} power, so convert with the duty
    cycle before choosing ($\S$8.2).
  \item Measurement questions are worded as ``describe how X and Y are measured''. X and Y
    are almost always \textbf{VSWR meter} and \textbf{some power sensor}, which is two
    separate answers glued together (72 Ash, 70 Bh).
\end{itemize}

\hr

%=======================================================================
\T{8.1 Microwave Measurements Vs Low-Frequency Measurements}

\Q{How are microwave measurements different from low-frequency measurements? List the major RF/MW measurement parameters}{\m{2}\m{3}\m{3}}
{\color{sub}\footnotesize \tF{4} \yr{\textbf{71 Bh}, \textbf{73 Bh}, \textbf{\texttt{79 Bh}}, \textbf{\texttt{81 Bh}}} \gap$\cdot$\gap always the opener to a longer measurement question}

\sbsr{0.52}{%
\lead{Why voltage and current stop being useful}
\begin{itemize}
  \item At low frequency the circuit is \textbf{lumped}: $V$ and $I$ are the same everywhere
    in a wire, so measure them and compute $Z$, power factor and phase angle.
  \item At microwaves the line is \textbf{distributed}: $V$ and $I$ are functions of
    \textbf{distance} along the line, so a ``the voltage'' does not exist. In a waveguide
    there is no unique voltage at all.
  \item \textbf{Power is the invariant}: on a lossless line the transmitted power is the same
    at every plane. So measure \textbf{power}, and get everything else from ratios of powers
    and from the standing-wave pattern.
  \item Probes and leads are a \textbf{significant fraction of a wavelength}, so any contact
    instrument loads and detunes what it measures. Measurement is made by
    \textbf{directional coupling} and by \textbf{sampling} the field, not by touching nodes.
  \item Circuit behaviour is described by \textbf{S-parameters}, which are themselves
    defined from travelling waves rather than terminal $V$ and $I$ ($\S$3.1).
  \item \textbf{Cost dodge used in every teaching laboratory}: rather than process microwave
    signals directly, the source is \textbf{1 kHz square-wave modulated}, the microwave
    signal is detected by a square-law crystal, and all the amplitude information is then
    read with an ordinary 1 kHz low-frequency amplifier, the \textbf{VSWR meter}.
\end{itemize}}{%
\lead{The measurables and their tools \mk{(81 Bh asks exactly this list)}}
\begin{tabularx}{\linewidth}{@{}L{2.35cm}X@{}}
\toprule
\hrow \thead{Measurable} & \thead{Instrument} \\
\midrule
Power (low) & bolometer (thermistor / barretter) bridge, Schottky diode sensor, thermocouple sensor \\
Power (medium, high) & calorimeter (static, circulating), calorimetric wattmeter, thermistor behind a calibrated attenuator or directional coupler \\
VSWR, $|\Gamma|$ & slotted line + VSWR meter, reflectometer \\
Impedance, unknown load & slotted line (minima-shift method), VNA, magic tee \\
Frequency & cavity (absorption) wavemeter, frequency counter, slotted line via $\lambda_g$ \\
Spectrum, harmonics & spectrum analyzer \\
S-parameters, insertion and return loss & network analyzer (VNA) \\
Attenuation, insertion loss & two directional couplers + VSWR meter \\
Noise figure & noise-source + Y-factor measurement ($\S$8.9) \\
Field strength, radiation zones & EMF survey meter with isotropic probe ($\S$8.10) \\
\bottomrule
\end{tabularx}}

\penalty0\vspace{2pt}

%=======================================================================
\T{8.2 Microwave Power: Bands, Average Vs Peak, Choosing the Sensor}

\Q{Choose a proper power measurement tool for the stated source and explain its working principle}{\m{8}\m{10}}
{\color{sub}\footnotesize \tF{4} \yr{\textbf{72 Ash}, 72 Ma, \textbf{73 Bh}, 74 Ma} \gap$\cdot$\gap 72 Ma 7.5 mW, 74 Ma airport surveillance radar, 73 Bh ``low microwave power''}

\begin{itemize}
  \item \textbf{Power} = energy dissipated or stored per unit time. At microwaves it is the
    \textbf{average} power that a sensor responds to:
    \[ P_{av} = \frac{1}{T}\int_0^T v(t)\,i(t)\,dt, \qquad
       P_{av} = P_{peak}\times \text{duty cycle}, \quad
       \text{duty} = \tau f_r = \tau/T \]
  \item Unit: \mk{$P(\mathrm{dBm}) = 10\log_{10}[P(\mathrm{mW})/1\,\mathrm{mW}]$}, so
    0 dBm = 1 mW, 30 dBm = 1 W, $-30$ dBm = 1 $\mu$W.
\end{itemize}

\par\penalty0\Needspace{13\baselineskip}
\lead{The decision table: this is the answer to every ``choose a tool'' question}
\begin{tabularx}{\linewidth}{@{}L{2.1cm}L{2.5cm}XL{2.9cm}@{}}
\toprule
\hrow \thead{Band} & \thead{Range} & \thead{Sensor, and why} & \thead{Section} \\
\midrule
\textbf{Low} & $<$ 10 mW & \textbf{resistance changes with power}: thermistor or barretter bolometer bridge (floor about $-20$ dBm), Schottky diode detector (down to $-70$ dBm = 100 pW), thermocouple sensor & $\S$8.3, $\S$8.4 \\
\textbf{Medium} & 10 mW to 10 W & thermistor bolometer behind a \textbf{calibrated attenuator or 20 dB directional coupler}; static or circulating calorimeter & $\S$8.3, $\S$8.5 \\
\textbf{High} & $>$ 10 W & \textbf{temperature changes with power}: calorimeter (static for moderate, circulating / flow for kW), calorimetric wattmeter & $\S$8.5 \\
\bottomrule
\end{tabularx}

\par\penalty0\Needspace{9\baselineskip}
\lead{Measurement procedure, common to every power meter (the ``zeroing'' marks)}
\begin{enumerate}[label=\arabic*., leftmargin=1.4em, itemsep=1pt]
  \item Attach a \textbf{calibrated sensor} to the line port where the unknown power appears;
    connect it to the power meter.
  \item \textbf{Turn the RF off and zero the meter} (``zero setting''): this removes sensor
    self-heating and ambient drift from the reading.
  \item Turn the RF on. The sensor reacts, the meter reads, and the reading is by
    \textbf{substitution}: the dc or low-frequency power displaced, or the temperature rise.
  \item Correct for the coupler or attenuator in front of the sensor, and for the
    \textbf{calibration factor} (mismatch plus sensor efficiency).
\end{enumerate}

\par\penalty0\Needspace{10\baselineskip}
\lead{The three PYQ cases worked \mk{(numbers checked in \code{meas.py})}}
\begin{itemize}
  \item \textbf{72 Ma, ``about 7.5 mW''} \m{10}: 7.5 mW $=$ 8.75 dBm, below 10 mW, so
    \textbf{low power}. Answer: \textbf{thermistor bolometer bridge}, sensitivity floor
    about $-20$ dBm $=$ 0.01 mW, so 7.5 mW sits comfortably inside its range. Then give the
    bridge working principle ($\S$8.3). A Schottky diode sensor is the alternative.
  \item \textbf{74 Ma, airport surveillance radar} \m{8}: a radar is quoted in \textbf{peak}
    power and is pulsed. Take a representative set, 1 MW peak, $\tau = 1\ \mu$s,
    PRF 1 kHz: duty $= 10^{-3}$, so $P_{av} = 1$ kW, which is \textbf{high power}. Answer:
    a \textbf{circulating (flow) calorimeter} on a dummy load, or a thermistor behind a
    high-directivity coupler plus a calibrated attenuator; report
    $P_{peak} = P_{av}/\text{duty}$. Never put a bolometer directly in the line: it would
    burn out.
  \item \textbf{73 Bh, ``low microwave power measurement device''} \m{8}: name the parameters
    first ($\S$8.1), then the \textbf{bolometer bridge} in full.
\end{itemize}

\penalty0\vspace{2pt}

%=======================================================================
\T{8.3 Bolometer Bridge Method}

\Q{Explain the working principle of bolometry used in microwave measurements}{\m{5}\m{6}\m{8}\m{10}}
{\color{sub}\footnotesize \tS{6} \yr{\textbf{70 Bh}, \textbf{72 Ash}, 72 Ma, \textbf{73 Bh}, \textbf{79 Ch}, \textbf{80 Ch}} \gap$\cdot$\gap 79 Ch for 10 marks, 80 Ch as a 5 mark short note}

\begin{itemize}
  \item A \textbf{bolometer} is a power sensor whose \textbf{resistance changes with
    temperature} as it absorbs microwave power. RF energy is turned to heat inside the
    element, the element warms, its resistance moves, and that resistance change is read.
  \item Bolometer impedances are \mk{100 to 200 $\Omega$}. The \textbf{mount} provides a good
    impedance match, low loss, thermal and mechanical isolation and shielding against
    leakage.
\end{itemize}

\par\penalty0\Needspace{12\baselineskip}
\sbsr{0.60}{%
\begin{itemize}
  \item \textbf{Two types}:
  \begin{itemize}
    \item \textbf{Barretter}: a short thin metal wire (platinum or tungsten, like a fuse),
      \textbf{positive} temperature coefficient. More delicate, so used only for very low
      power, a few mW at most.
    \item \textbf{Thermistor}: a semiconductor bead, \textbf{negative} temperature
      coefficient. Small, rugged, easily mounted in a waveguide or coaxial mount; the usual
      choice. Sensitivity floor about $-20$ dBm.
  \end{itemize}
\end{itemize}}{%
\figT{c8_barretter.png}
{\color{sub}\scriptsize Resistance vs temperature: thermistor falls (negative coefficient),
barretter rises (positive). Karkee, Ch-8 deck slide 8.}}

\par\penalty0\Needspace{11\baselineskip}
\sbsr{0.60}{%
\begin{itemize}
  \item \textbf{Principle}: the element sits in one arm of a \textbf{Wheatstone bridge}.
    With no RF, $R_5$ is adjusted to balance the bridge. Incident microwave power heats the
    element, its resistance changes, the bridge unbalances, and the galvanometer or voltmeter
    current is \textbf{calibrated directly in power}.
  \item \textbf{Substitution}: because the element responds only to total heating, the RF
    power is exactly the \textbf{dc power that has to be withdrawn} to restore balance. This
    is what makes it an absolute measurement rather than a comparison.
\end{itemize}}{%
\figT{c8_bridge.png}
{\color{sub}\scriptsize Bolometer mount in one arm of a Wheatstone bridge; the galvanometer
current measures the incident power. Deck slide 9.}}

\par\penalty0\Needspace{16\baselineskip}
\creamq{Single bridge, its two faults, and the double-channel (double bridge) cure}{\m{2}\m{4}\m{8}}
{\color{sub}\footnotesize \tF{3} \yr{\texttt{80 Ba}, \textbf{\texttt{80 Bh}}, \texttt{81 Ba}} \gap$\cdot$\gap 80 Ba asks only the \textbf{limitations} \m{2}; 81 Ba is a 4 mark short note}

\sbsr{0.50}{%
\lead{Single bridge}
\begin{itemize}
  \item Balanced under \textbf{zero incident power} by adjusting $R_5$; the bolometer arm is
    $R_4$.
  \item Microwave power into $R_4$ changes its resistance, unbalances the bridge, and the
    non-zero output is read on a voltmeter calibrated in power.
  \item With dc bias voltage $E_1$ at balance and $E_2$ after RF is applied, the change
    $(E_1 - E_2)$ is directly proportional to the microwave power.
  \item \mk{Two faults} (this is the whole of 80 Ba's 2 marks):
  \begin{enumerate}[label=\arabic*., leftmargin=1.4em]
    \item A \textbf{mismatch} at the microwave input reflects power, so the element absorbs
      less than is incident: the reading is wrong and there is no way to tell.
    \item The thermistor is \textbf{sensitive to ambient temperature}: room drift is read as
      RF power. A false reading with no warning.
  \end{enumerate}
\end{itemize}}{%
\figT{c8_singlebridge.png}
{\color{sub}\scriptsize Self-balancing single bolometer bridge. Deck slide 10.}}

\par\penalty0\Needspace{18\baselineskip}
\sbsr{0.52}{%
\lead{Double-channel bolometer bridge}
\begin{itemize}
  \item \textbf{Two identical bridges}. The \textbf{upper} one carries the RF sensing
    thermistor and measures the power. The \textbf{lower} one carries a matched thermistor
    with \textbf{no RF}, exposed to the same ambient, and compensates temperature drift by
    forcing $V_1 = V_2$.
  \item The extra heating caused by mismatch is taken out automatically by \textbf{negative
    dc feedback}: the self-balancing loop lowers the dc power in the sensing thermistor until
    balance is restored, that is until the \textbf{net resistance change is zero}.
  \item \textbf{Zero setting}: with no input, adjust $E_1 = E_2 = E_0$. Without RF the dc
    voltage across the sensor at balance is $E_1/2$; with RF it is $E_2/2$.
  \item Average input power is the dc power the RF displaced:
    \[ P_{av} = \frac{E_1^2}{4R} - \frac{E_2^2}{4R}
              = \frac{(E_1-E_2)(E_1+E_2)}{4R} \]
  \item If ambient drift moves both rails by $\Delta E$,
    \[ P_{av} + \Delta P = \frac{(E_1-E_2)(E_1-E_2+2\Delta E)}{4R} \]
    and since $E_1 + E_2 \gg \Delta E$ in practice, $\Delta P \approx 0$: the meter reads the
    line above. \textbf{That inequality is why the second bridge works.}
\end{itemize}}{%
\figT{c8_doublebridge.png}
{\color{sub}\scriptsize Double-channel bolometer bridge: upper bridge measures, lower bridge
compensates ambient temperature. Deck slide 12.}}

\par\penalty0\Needspace{8\baselineskip}
\begin{itemize}
  \item \textbf{Worked, to fix the formula} (\code{meas.py}): $R = 200\ \Omega$,
    $E_1 = 6$ V, $E_2 = 4$ V. Sensor dc power falls from 45 mW to 20 mW, so
    $P_{av} = (36-16)/800 = $ \mk{25 mW}. A $+50$ mV ambient drift on both rails would move a
    \textbf{single} bridge to 25.25 mW, an error of 1\%, which the second bridge removes.
\end{itemize}

\penalty0\vspace{2pt}

%=======================================================================
\T{8.4 Schottky Diode Detector and Thermocouple Sensor}

\Q{Other low-power sensors: the square-law diode detector and the thermocouple}{\pill{gdL}{gd}{syllabus 8.4}}
{\color{sub}\footnotesize no question names them, but both are legitimate answers to ``a low power measurement device'' (73 Bh, 72 Ash, 72 Ma)}

\sbsr{0.60}{%
\lead{Schottky barrier diode detector}
\begin{itemize}
  \item A zero-biased \textbf{metal-semiconductor (point contact) diode} used as a
    \textbf{square-law detector}: output is proportional to the \textbf{square} of the input
    voltage, therefore to input \textbf{power}.
  \item RF input is applied to $R_1$ and passes through $R_2$; the diode rectifies, and the
    resulting low-frequency or dc current is read.
  \item Sensitivity down to \mk{$-70$ dBm (100 pW)}, up to about 18 GHz. It is the fastest
    and most sensitive of the low-power sensors.
  \item The circuit is arranged so the input match does not depend on the diode resistance,
    which is a strong function of temperature.
  \item \textbf{Caution}: it is square-law only over a limited range. Above roughly $-20$ dBm
    it enters a transition region and eventually becomes \textbf{linear}, which is exactly
    why the VSWR meter misreads a high SWR ($\S$8.6).
\end{itemize}}{%
\figT{c8_schottky.png}
{\color{sub}\scriptsize Schottky barrier diode detector. Deck slide 6.}}

\par\penalty0\Needspace{11\baselineskip}
\sbsr{0.60}{%
\lead{Thermocouple sensor}
\begin{itemize}
  \item A \textbf{thermocouple} is a junction of two dissimilar metals or semiconductors
    (here n-type Si). A thin film of \textbf{tantalum-nitride} resistive load is deposited on
    a Si substrate and forms one electrode.
  \item Microwaves are absorbed in the resistive load, heat one end of the junction more than
    the other, and the junction generates an \textbf{emf proportional to incident power}.
  \item $C_2$ is the RF bypass capacitor, $C_1$ the input coupling capacitor (dc block). Emfs
    of the parallel thermocouples add across $C_2$; the output leads are at RF ground, so the
    dc voltmeter reads pure dc.
  \item For a square-wave modulated signal the peak power follows from the average:
    $P_{peak} = P_{av}\,T/\tau$, $T$ the period and $\tau$ the pulse width.
    \mk{Same relation as the radar case in $\S$8.2.}
  \item Wider dynamic range and better match stability than a bolometer, and no dc bias, but
    no substitution: it must be calibrated.
\end{itemize}}{%
\figT{c8_thermocouple.png}
{\color{sub}\scriptsize Thermocouple power sensor. Deck slide 14.}}

\penalty0\vspace{2pt}

%=======================================================================
\T{8.5 Calorimetric Methods}

\Q{What is calorimetry in microwave? Explain how microwave power is measured with a static calorimeter}{\m{2}\m{5}\m{6}\m{7}\m{10}}
{\color{sub}\footnotesize \tS{6} \yr{70 Ma, 71 Ma, \textbf{71 Bh}, \textbf{76 Bh}, \textbf{77 Ch}, \textbf{\texttt{81 Bh}}} \gap$\cdot$\gap 76 Bh says ``static and dry'', which is the same instrument: a dry dielectric load, no flowing water}

\lead{Calorimetry: the 2 mark definition}
\begin{itemize}
  \item \textbf{Calorimetry} measures microwave power by \textbf{converting it entirely into
    heat}, absorbing that heat in a known medium, and measuring the medium's
    \textbf{temperature rise}. It is the standard method for \textbf{medium and high power},
    and it is an \textbf{absolute} measurement: nothing is compared against a reference
    source, only against the specific heat of a fluid.
  \item \textbf{Two heating arrangements} (both used with either calorimeter below):
  \begin{itemize}
    \item \textbf{Direct heating}: the rate of heat production is found from the temperature
      rise of the \textbf{dissipating medium itself}.
    \item \textbf{Indirect heating}: heat is \textbf{transferred to another medium} before
      it is measured.
  \end{itemize}
\end{itemize}

\par\penalty0\Needspace{13\baselineskip}
\sbsr{0.60}{%
\lead{Static (dry) calorimeter}
\begin{itemize}
  \item A \textbf{50 $\Omega$ coaxial line terminated in a lossy dielectric load} with a high
    hysteresis loss, thermally isolated from its surroundings. No fluid circulates, hence
    ``\textbf{dry}''.
  \item The load dissipates all the microwave power; a thermometer or thermocouple reads the
    rise $T$ of the thermometric medium over a time $t$:
    \[ \boxed{P = \frac{4.187\, m\, C_p\, T}{t}\ \ \mathrm{watts}} \]
    $m$ mass in grams, $C_p$ specific heat in cal/g, $T$ rise in $^\circ$C, $t$ in seconds,
    and \mk{4.187 is joules per calorie}.
  \item \textbf{Where the formula comes from}: $m C_p T$ is the heat in \textbf{calories},
    dividing by $t$ makes it calories per second, and 4.187 converts to watts. Quote this and
    the formula cannot be mis-remembered.
  \item \textbf{Worked} (\code{meas.py}): 100 g of water raised 5$^\circ$C in 60 s is
    $4.187\times100\times1\times5/60 = $ \mk{34.89 W}.
  \item \textbf{Limits}: slow (thermal inertia), and the load must stay isolated, so it suits
    steady medium power rather than a kilowatt transmitter.
\end{itemize}}{%
\figT{c8_static.png}
{\color{sub}\scriptsize Static calorimeter: RF into an isolated calorimetric body, compared
against low-frequency power. Deck slide 18.}}

\par\penalty0\Needspace{14\baselineskip}
\creamq{Circulating (flow) calorimeter; circulating Vs flow calorimetry}{\m{4}\m{5}\m{8}}
{\color{sub}\footnotesize \tF{4} \yr{70 Ma, \textbf{75 Bh}, \textbf{\texttt{81 Bh}}, \texttt{82 Ba}} \gap$\cdot$\gap 82 Ba is a 4 mark short note; 81 Bh pairs it with the measurables list of $\S$8.1}

\sbsr{0.58}{%
\begin{itemize}
  \item The calorimeter fluid (usually \textbf{water}) is \textbf{pumped continuously}
    through a water load in the waveguide. Microwave power heats it in transit, so the
    \textbf{outlet} temperature $T_2$ exceeds the \textbf{inlet} $T_1$:
    \[ \boxed{P = 4.187\, v\, d\, C_p\,(T_2 - T_1)\ \ \mathrm{watts}} \]
    $v$ flow rate in cc/s, $d$ specific gravity in g/cc, $C_p$ specific heat in cal/g.
  \item \textbf{Same law as the static one}: $v\,d$ is grams per second, which is the static
    formula's $m/t$. Confirmed in \code{meas.py}.
  \item \textbf{Worked}: 10 cc/s of water with a 2$^\circ$C rise is \mk{83.7 W}. To take
    1 kW with a 10$^\circ$C rise needs \mk{23.9 cc/s}.
  \item \textbf{Instrumentation}: pump, flow meter, glass tube through the guide, inlet and
    outlet thermocouples, heat exchanger to return the water cool.
\end{itemize}}{%
\figT{c8_circulating.png}
{\color{sub}\scriptsize Circulating (flow) calorimeter: water load, pump, flow meter, inlet
and outlet thermometers. Deck slide 19.}}

\par\penalty0\Needspace{11\baselineskip}
\lead{Static Vs circulating (the ``differentiate'' answer)}
\begin{tabularx}{\linewidth}{@{}L{2.4cm}XX@{}}
\toprule
\hrow & \thead{Static (dry)} & \thead{Circulating (flow)} \\
\midrule
Load & fixed dry dielectric, thermally isolated & water pumped through continuously \\
Reading & temperature rise per unit \textbf{time} & \textbf{steady} inlet to outlet difference \\
Formula & $4.187\,mC_pT/t$ & $4.187\,vdC_p(T_2{-}T_1)$ \\
Power handled & medium; the load stores heat and saturates & \textbf{high}, kW upward: the flow carries heat away \\
Speed & slow, must wait for a measurable rise & continuous once flow is steady \\
Errors & heat leakage, isolation, thermal inertia & flow-rate stability, thermometer placement, thermal inertia \\
\bottomrule
\end{tabularx}
\begin{itemize}
  \item \textbf{On the wording ``circulating Vs flow''} (70 Ma) \added{}: the sources use the
    two words for the \textbf{same} instrument. Answer the question that is meant, which is
    \textbf{static vs circulating}, and say so in one line: if a distinction is forced,
    \textbf{circulating} returns the same fluid round a closed loop through a heat exchanger,
    while a plain \textbf{flow} calorimeter passes fluid once and discards it.
\end{itemize}

\par\penalty0\Needspace{16\baselineskip}
\creamq{Power measurement using a calorimeter wattmeter}{\m{6}}
{\color{sub}\footnotesize \tP{1} \yr{\texttt{80 Ba}} \gap$\cdot$\gap paired with the single-bridge limitations of $\S$8.3}

\sbsr{0.56}{%
\begin{itemize}
  \item A calorimeter run as a \textbf{self-balancing comparison bridge}, so the answer comes
    out as a meter reading instead of a temperature calculation.
  \item The unknown RF power is checked against a \textbf{1200 cps (Hz) comparison power} in
    a bridge circuit. Two \textbf{temperature-sensitive resistors} serve as gauges.
  \item \textbf{Operation}:
  \begin{enumerate}[label=\arabic*., leftmargin=1.4em]
    \item The unknown RF heats an \textbf{input load resistor}.
    \item That resistor and one gauge sit in close thermal contact, so the gauge warms and
      \textbf{unbalances the bridge}.
    \item The unbalance signal is \textbf{amplified} and applied to the \textbf{comparison
      load resistor}, which is next to the second gauge, and \textbf{rebalances} the bridge.
    \item The meter reads the \textbf{power supplied to the comparison load}, which at
      balance equals the unknown RF power.
  \end{enumerate}
  \item Both loads and gauges are \textbf{immersed in an oil stream} for efficient, equal
    heat transfer.
  \item \textbf{Why it is better than a bare calorimeter}: it is a null method, so gauge
    calibration and ambient drift cancel, and it reads continuously.
\end{itemize}}{%
\figT{c8_wattmeter.png}
{\color{sub}\scriptsize Calorimetric wattmeter: input load and comparison load, two
temperature gauges, 1200 cps comparison power, oil stream. Deck slide 20.}}

\penalty0\vspace{2pt}

%=======================================================================
\T{8.6 VSWR Measurement: Slotted Line and VSWR Meter}

\Q{Describe how standing waves and microwave power are measured with a VSWR meter}{\m{2}\m{4}}
{\color{sub}\footnotesize \tF{3} \yr{\textbf{70 Bh}, \textbf{72 Ash}, \textbf{\texttt{79 Bh}}} \gap$\cdot$\gap always part A, with a power sensor as part B: 72 Ash ``low power measurement'' \m{8}, 70 Bh ``bolometry'' \m{6}}

\sbsr{0.58}{%
\lead{Slotted line carriage: where the standing wave is seen}
\begin{itemize}
  \item A \textbf{coaxial E-field probe} penetrates a rectangular waveguide (or coaxial line)
    through a \textbf{longitudinal slot} cut in the centre of the broad wall, over 2 to 3
    wavelengths.
  \item The slot lies where the \textbf{wall current has no transverse component}, so cutting
    it does not disturb the field. Both ends are \textbf{tapered} to zero width to reduce
    discontinuity, and the probe is thin and shallow for the same reason.
\end{itemize}}{%
\figT{c8_slotted.png}
{\color{sub}\scriptsize Slotted-line carriage: longitudinal slot on the broad wall, probe,
crystal detector, vernier scale. Deck slide 22.}}

\par\penalty0\Needspace{10\baselineskip}
\sbsr{0.58}{%
\begin{itemize}
  \item The carriage moves the probe along the slot at constant depth, so coupling is uniform,
    and a \textbf{vernier scale} reads its position. It carries a tunable stub and a crystal
    detector, whose 1 kHz output goes to the VSWR meter.
  \item \textbf{It measures}: standing-wave pattern and \textbf{VSWR}, \textbf{wavelength}
    ($\lambda_g = 2\times$ distance between successive minima), and by the minima-shift
    method \textbf{impedance, reflection coefficient and return loss} ($\S$8.7).
\end{itemize}}{%
\figT{c8_vswrsetup.png}
{\color{sub}\scriptsize The pattern the probe traces: maxima and minima $\lambda_g/2$ apart,
each minimum $\lambda_g/4$ from its neighbouring maximum. Deck slide 27.}}

\par\penalty0\Needspace{12\baselineskip}
\sbsr{0.54}{%
\lead{VSWR meter}
\begin{itemize}
  \item A \textbf{high-gain, high-$Q$, low-noise voltage amplifier} tuned to a \textbf{fixed
    1 kHz}, the frequency at which the microwave source is square-wave modulated. Overall
    gain about 125 dB, switchable in 10 dB steps.
  \item Input is the \textbf{detected} output of the crystal detector; output goes to a
    \textbf{square-law calibrated} voltmeter, which therefore reads $V_{max}/V_{min}$
    directly once it has been set to 1.0 at a maximum.
  \item \textbf{Procedure} (low VSWR):
  \begin{enumerate}[label=\arabic*., leftmargin=1.4em]
    \item Set the attenuator to about 10 dB, tune the probe stub for maximum reading.
    \item Slide the carriage to a \textbf{maximum} and adjust the \textbf{gain} until the
      meter reads \textbf{1.0} (0 dB).
    \item Move to the adjacent \textbf{minimum}. The meter now reads \mk{$S$ directly}.
  \end{enumerate}
  \item \textbf{Three scales}: SWR NORMAL top (1 to 4), SWR NORMAL bottom (3.2 to 10),
    EXPANDED (1 to 1.3) for a nearly matched load, plus a dB scale. For $S$ between 10 and
    40 increase sensitivity by 20 dB, read on the 1 to 4 scale and \textbf{multiply by 10}.
\end{itemize}}{%
\figT{c8_vswrmeter.png}
{\color{sub}\scriptsize 1 kHz modulation, mixer, amplifier, detector: the meter reads
$V_{max}/V_{min}$. Deck slide 31.}
\vspace{2pt}
\figT{c8_lowvswr.png}
{\color{sub}\scriptsize Low-VSWR bench: source, attenuator, slotted line, crystal detector,
power meter. Deck slide 28.}}

\par\penalty0\Needspace{18\baselineskip}
\creamq{How is the VSWR of a microwave transmitter measured when VSWR $>$ 10?}{\m{6}}
{\color{sub}\footnotesize \tP{1} \yr{\textbf{\texttt{79 Bh}}} \gap$\cdot$\gap \mk{the double-minimum method}: the one question in the chapter with a derivation in it}

\sbsr{0.56}{%
\begin{itemize}
  \item \textbf{Why the direct reading fails.} With $S > 10$ the gap between $V_{max}$ and
    $V_{min}$ is large. To read $V_{max}$ the detector is driven out of its
    \textbf{square-law} region (diode current past about 20 mA, response turning linear),
    while at $V_{min}$ it is in the noise. A meter calibrated on $I = kV^2$ is then simply
    wrong, and the error grows with $S$.
  \item \textbf{The cure}: never look at the maximum. Take \textbf{two readings either side
    of a single minimum}, where the level is low and the detector is safely square-law. This
    is the \textbf{double-minimum method}.
  \item \textbf{Procedure}:
  \begin{enumerate}[label=\arabic*., leftmargin=1.4em]
    \item Terminate the line with the load and find a \textbf{voltage minimum} on the slotted
      line; note its position and meter reading.
    \item Move the probe either side until the detected \textbf{power is twice the minimum}
      (voltage $\sqrt2\,V_{min}$). Record the two positions $d_1$ and $d_2$;
      $\Delta x = d_2 - d_1$.
    \item Short the line and measure $\lambda_g$ as \textbf{twice} the distance between two
      successive minima.
    \item Compute $S$.
  \end{enumerate}
\end{itemize}}{%
\figT{c8_highvswr.png}
{\color{sub}\scriptsize The two ``twice minimum power'' points either side of a minimum,
$\Delta x = d_2 - d_1$. Deck slide 29.}}

\begin{itemize}
  \item \textbf{The formula, and where it comes from.} Near a minimum the detected quantity
    is $|V|^2 = V_{min}^2 + (V_{max}^2 - V_{min}^2)\sin^2\beta\delta$ with $\delta$ measured
    from the minimum. Setting $|V|^2 = 2V_{min}^2$ gives $(S^2-1)\sin^2\beta\delta = 1$, so
    with $\Delta x = 2\delta$ and $\beta = 2\pi/\lambda_g$:
    \[ \boxed{S = \sqrt{1 + \frac{1}{\sin^2(\pi\Delta x/\lambda_g)}}}
       \qquad\longrightarrow\qquad S \approx \frac{\lambda_g}{\pi\,\Delta x}
       \ \ \text{for large } S \]
  \item \textbf{Both forms verified} in \code{meas.py} against a simulated standing wave: at
    $S = 25$ the exact form returns 25.00 and the short form 24.97 (0.12\% error); at
    $S = 60$, 0.04\%. \mk{The short form is only for large $S$}: at $S = 2$ it returns 1.63,
    19\% out, which is precisely why the direct method is used there instead.
  \item \textbf{Threshold}: the deck sets the changeover at $S > 10$ (which is the wording of
    79 Bh), Das at $S > 20$. Say ``above about 10, where the detector leaves its square-law
    region'' and the distinction never matters.
\end{itemize}

\penalty0\vspace{2pt}

%=======================================================================
\T{8.7 Impedance, Reflection Coefficient and Frequency}

\Q{Measurement of an unknown load, of reflection coefficient, and of frequency}{\pill{gdL}{gd}{syllabus 8.5--8.8}}
{\color{sub}\footnotesize no direct PYQ, but it is four of the nine syllabus sub-topics, and it is the reasoning behind 74 Bh's ``antenna as DUT'' \m{8} in $\S$7.8}

\sbsr{0.53}{%
\lead{Unknown impedance from the slotted line (minima-shift method)}
\begin{itemize}
  \item Impedance is complex, so \textbf{two} numbers are needed: the slotted line supplies
    \textbf{$S$} (giving $|\Gamma|$) and \textbf{$d_{min}$} (giving the phase).
  \item \textbf{Set-up 1}: load connected, measure $S$ and the position of the first minimum.
    \textbf{Set-up 2}: replace the load with a \textbf{short}, which puts a minimum at the
    load plane, and measure the \textbf{shift}.
  \item Shift \textbf{towards the generator} means the load is \textbf{capacitive}; shift
    \textbf{towards the load} means \textbf{inductive}.
  \item The relations:
    \[ \rho = \frac{S-1}{S+1},\quad \phi_L = 2\beta d_{min} - \pi,\quad
       \beta = \frac{2\pi}{\lambda_g},\quad Z_L = Z_0\frac{1+\Gamma}{1-\Gamma} \]
    with $\lambda_g = 2\times$(distance between successive minima), $\Gamma=\rho e^{j\phi_L}$.
  \item Equivalent closed form, useful when no chart is at hand, $t = \tan\beta d_{min}$:
    \[ Z_L = Z_0\,\frac{1 - jSt}{S - jt} \]
  \item \textbf{Both routes checked} in \code{meas.py}: a simulated pattern from
    $Z_L = 60 - j80$ on a 50 $\Omega$ line is inverted back to $60 - j80$ by each.
  \item On the Smith chart: draw the $S$ circle, start at the short-circuit point and move
    \textbf{towards the load} by $d_{min}/\lambda_g$; the intersection is $z_L$.
\end{itemize}}{%
\figT{c8_impedance.png}
{\color{sub}\scriptsize The two set-ups: unknown load, then shorted termination. Deck
slide 24.}
\vspace{2pt}
\figT{c8_reflectometer.png}
{\color{sub}\scriptsize Reflectometer: two 20 dB directional couplers sample forward and
reverse waves. Deck slide 26.}}

\par\penalty0\Needspace{11\baselineskip}
\sbsr{0.53}{%
\lead{Reflectometer: $|\Gamma|$ without a slotted line}
\begin{itemize}
  \item \textbf{Two directional couplers} face opposite ways: the first samples the
    \textbf{forward} wave at port 4, the second the \textbf{reverse} wave at port 3.
  \item With coupling $C \ll 1$ the ratio of the two detected outputs is the reflection
    coefficient itself:
    \[ \left|\frac{b_3}{b_4}\right| = \sqrt{1-C^2}\,|\Gamma_L| \approx |\Gamma_L| \]
  \item Calibrate by shorting port 2 ($|\Gamma| = 1$) and setting the meter to unity, then
    connect the load. Everything else follows:
    \[ \mathrm{VSWR} = \frac{1+|\Gamma_L|}{1-|\Gamma_L|}, \qquad
       \mathrm{RL(dB)} = -20\log_{10}|\Gamma_L| \]
  \item \textbf{Broadband and fast}, but accuracy is set by coupler \textbf{directivity}, so
    it suits \textbf{low} VSWR; a high SWR still wants the slotted line.
  \item Related loss family, all measured the same way:
    \textbf{insertion loss} $=$ reflection loss $+$ attenuation loss, with reflection loss
    $=-10\log(1-|\Gamma|^2) = -10\log[4S/(1+S)^2]$. \mk{Return loss compares $P_r$ with
    $P_i$; insertion loss compares $P_0$ with $P_i$} ($\S$3.4).
\end{itemize}}{%
\lead{Frequency measurement}
\begin{itemize}
  \item \textbf{Cavity (absorption) wavemeter}: a cylindrical cavity with a \textbf{variable
    short}, coupled to the guide through a hole in the narrow wall. Tuning the plunger
    changes the resonant frequency; at resonance the cavity \textbf{absorbs}, and the power
    meter downstream shows a \textbf{dip}. The plunger scale is calibrated in frequency.
    \begin{itemize}
      \item $TE_{011}$ has the highest $Q$ and no axial current, but it is a higher-order
        mode, so practical wavemeters use the dominant \textbf{$TM_{010}$}, excited by
        magnetic coupling. Polytron absorber behind the plunger stops oscillation above it.
      \item Accuracy follows $Q$: \mk{1\% at $Q = 1000$, 0.005\% at $Q = 50\,000$}.
    \end{itemize}
  \item \textbf{Slotted line}: short the guide, measure $\lambda_g$ as twice the minima
    spacing, then
    \[ \frac{1}{\lambda^2} = \frac{1}{\lambda_g^2} + \frac{1}{\lambda_c^2},
       \qquad f = \frac{c}{\lambda},\qquad \lambda_c = 2a\ (TE_{10}) \]
    Worked (\code{meas.py}): WR-90 ($a = 2.286$ cm) with $\lambda_g = 4.00$ cm gives
    \mk{9.96 GHz}.
  \item \textbf{Frequency counter} with a \textbf{down-conversion} front end: the unknown is
    mixed with a selected harmonic $nf_1$ of a known oscillator, and the counter reads
    $f_x = $ counter $+ nf_1$. Direct and most accurate.
\end{itemize}}

\penalty0\vspace{2pt}

%=======================================================================
\T{8.8 Spectrum Analyzer}

\Q{Spectrum analyzer (short note): working principle and block diagram}{\m{4}\m{5}}
{\color{sub}\footnotesize \tP{2} \yr{\textbf{\texttt{80 Bh}}, \textbf{80 Ch}} \gap$\cdot$\gap short note both times, so a block diagram plus six lines is the full answer}

\sbsr{0.56}{%
\begin{itemize}
  \item A \textbf{swept-tuned superheterodyne receiver} that plots \textbf{amplitude against
    frequency}: frequency on the horizontal axis, amplitude on the vertical. It shows what an
    oscilloscope cannot, namely the \textbf{signal spectrum}.
  \item \textbf{Working}:
  \begin{enumerate}[label=\arabic*., leftmargin=1.4em]
    \item The input microwave signal is mixed with a \textbf{local oscillator swept
      electronically} between two limits at a linear rate by a \textbf{sawtooth} sweep
      voltage.
    \item The difference product is amplified by a \textbf{narrow-bandwidth IF amplifier},
      then detected and \textbf{video amplified}.
    \item The same sawtooth drives the horizontal deflection, with \textbf{zero flyback
      time}, so horizontal position tracks the LO frequency while the detected amplitude
      gives vertical deflection.
  \end{enumerate}
\end{itemize}}{%
\figT{c8_spectrum.png}
{\color{sub}\scriptsize Block diagram: mixer, swept local oscillator, narrow-band IF
amplifier, detector, video amplifier, display. Deck slide 33.}}

\begin{itemize}
  \item \textbf{Measures}: carrier power, harmonics, spurious products, sidebands, phase
    noise, occupied bandwidth; on a modern analyzer, demodulated WLAN, GSM or LTE signals.
  \item \textbf{Four design parameters}: sweep \textbf{rate}, sweep \textbf{range}, IF
    \textbf{bandwidth}, IF \textbf{centre frequency}.
  \item \textbf{Resolution is the IF bandwidth (RBW)}. Narrow RBW resolves closer tones but
    needs a \textbf{slower sweep}, because the IF filter must have time to settle: sweep time
    rises as $1/\mathrm{RBW}^2$. Keep the span as small as the measurement allows.
  \item \textbf{Image response}: any signal $f_i = f_0 \pm f_{IF} = f_s \pm 2f_{IF}$ also
    beats to the IF. Choose $f_{IF}$ \textbf{high} to push the image out of band. Worked:
    $f_{IF} = 450$ kHz puts the image \mk{900 kHz} from the signal; $f_{IF} = 2$ MHz moves it
    to \mk{4 MHz} away, easily filtered (\code{meas.py}).
  \item \textbf{Types}: swept-tuned (superheterodyne), FFT, real-time, audio.
\end{itemize}

\penalty0\vspace{2pt}

%=======================================================================
\T{8.9 Network Analyzer, VNA and the Two-Port Model}

\Q{Describe the working principle of a network analyzer; VNA (short note)}{\m{5}\m{8}}
{\color{sub}\footnotesize \tF{3} \yr{\textbf{71 Bh}, 73 Ma, \textbf{78 Ch}} \gap$\cdot$\gap 78 Ch as a 5 mark short note on ``VNA type network analyzer''; 71 Bh asks the two-port model below}

\sbsr{0.52}{%
\lead{Why it exists, and how it works}
\begin{itemize}
  \item The slotted line measures amplitude and phase \textbf{one frequency at a time}, so
    broadband testing with it is slow and expensive. A \textbf{network analyzer} measures
    both \textbf{over a wide band} in a reasonable time.
  \item \textbf{Principle: everything is measured against a reference.}
  \begin{enumerate}[label=\arabic*., leftmargin=1.4em]
    \item A \textbf{sweep oscillator} feeds a \textbf{power divider}, splitting the signal
      into a \textbf{test} channel and a \textbf{reference} channel.
    \item The test signal passes through the \textbf{DUT}; the reference passes through a
      \textbf{phase-equalising length} of line.
    \item Microwave frequencies cannot be processed directly, so a \textbf{harmonic frequency
      converter} brings both to a fixed IF. A \textbf{phase-locked loop} makes the local
      oscillator track the reference channel, which is what allows a \textbf{swept}
      measurement.
    \item Conversion is in two steps, RF to an MHz IF, then to a kHz IF, where the
      \textbf{magnitude and phase comparators} compare test against reference.
  \end{enumerate}
  \item \textbf{Reflection and transmission} are separated by a test unit built from
    \textbf{dual directional couplers}: $S_{11}$ and $S_{21}$ from one direction, then source
    and load are interchanged for $S_{22}$ and $S_{12}$. Each $S_{ij}$ is a
    \textbf{ratio of a coupler output to the reference}, with the phase as the difference
    $\phi_j - \phi_i$.
  \item \textbf{Two types}: \textbf{SNA} (scalar) measures amplitude only; \textbf{VNA}
    (vector) measures \textbf{amplitude and phase}, so it can apply \textbf{vector error
    correction} and is far more accurate.
\end{itemize}}{%
\figT{c8_vna.png}
{\color{sub}\scriptsize VNA: DUT alters the reference in magnitude (resistive) and phase
(reactive); two comparators output the magnitude and phase vectors. Deck slide 36.}
\vspace{2pt}
\begin{tabularx}{\linewidth}{@{}L{1.55cm}XX@{}}
\toprule
\hrow & \thead{Spectrum an.} & \thead{Network an.} \\
\midrule
Measures & unknown signals: power, harmonics, sidebands, phase noise & known stimulus: $\Gamma$, $S$-parameters, insertion and return loss \\
Accuracy & lower & higher, vector error correction \\
IF BW & wider & narrower \\
Quantity & scalar only & amplitude \textbf{and} phase \\
Sweep & frequency & frequency \textbf{and} power \\
\bottomrule
\end{tabularx}}

\par\penalty0\Needspace{14\baselineskip}
\creamq{Design a two-port network model and derive the required parameters}{\m{10}}
{\color{sub}\footnotesize \tP{1} \yr{\textbf{71 Bh}} \gap$\cdot$\gap the only 10 mark question in $\S$8.9; it wants the \textbf{set} of two-port descriptions and why $S$ is the microwave one}

\begin{itemize}
  \item \textbf{The model}: a black box with two ports, port 1 $(V_1, I_1)$ and port 2
    $(V_2, I_2)$, both currents taken as flowing \textbf{into} the box. Any two of the four
    are independent, which is what gives the six standard parameter sets.
\end{itemize}
\begin{tabularx}{\linewidth}{@{}L{1.5cm}L{4.4cm}XL{3.2cm}@{}}
\toprule
\hrow \thead{Set} & \thead{Defining equations} & \thead{Element found by} & \thead{Needs} \\
\midrule
$Z$ (open-circuit) & $V_1 = Z_{11}I_1 + Z_{12}I_2$, $V_2 = Z_{21}I_1 + Z_{22}I_2$ & $Z_{11} = V_1/I_1|_{I_2=0}$ & \textbf{open} circuit at a port \\
$Y$ (short-circuit) & $I_1 = Y_{11}V_1 + Y_{12}V_2$, $I_2 = Y_{21}V_1 + Y_{22}V_2$ & $Y_{11} = I_1/V_1|_{V_2=0}$ & \textbf{short} circuit at a port \\
$h$ (hybrid) & $V_1 = h_{11}I_1 + h_{12}V_2$, $I_2 = h_{21}I_1 + h_{22}V_2$ & $h_{21} = I_2/I_1|_{V_2=0}$ & one short, one open \\
$ABCD$ (chain) & $V_1 = AV_2 - BI_2$, $I_1 = CV_2 - DI_2$ & cascades by matrix product & short and open \\
$S$ (scattering) & $b_1 = S_{11}a_1 + S_{12}a_2$, $b_2 = S_{21}a_1 + S_{22}a_2$ & $S_{11} = b_1/a_1|_{a_2=0}$ & \textbf{matched} $Z_0$ termination \\
\bottomrule
\end{tabularx}
\begin{itemize}
  \item \textbf{Why $S$ at microwaves} (the same argument as $\S$3.1, and the reason this
    question sits in the measurements chapter): $Z$, $Y$, $h$ and $ABCD$ all need a
    \textbf{short or open circuit} at a port. At microwaves a short is not a short (lead
    inductance), an open radiates, and both are \textbf{frequency dependent} and may
    destabilise an active device. A \textbf{matched load} is realisable at every frequency,
    and $a$, $b$ are travelling-wave quantities that a directional coupler can sample: that
    is exactly what the network analyzer above does, which is why $S$-parameters are the set
    that can actually be \textbf{measured}.
  \item \textbf{Reciprocity} $S_{12} = S_{21}$, \textbf{losslessness} $[S]^\dagger[S] = [I]$,
    \textbf{matched} $S_{ii} = 0$; conversions to and from $Z$ exist,
    $[S] = ([Z]-[I])([Z]+[I])^{-1}$ for $Z$ normalised to $Z_0$ ($\S$3.2).
\end{itemize}

\par\penalty0\Needspace{8\baselineskip}
\lead{Noise figure measurement \pill{gdL}{gd}{syllabus 8.9, never asked}}
\begin{itemize}
  \item \textbf{Y-factor method}: drive the DUT with a calibrated \textbf{noise source} that
    switches between hot and cold. $Y$ is the ratio of the two output powers, and with the
    source's \textbf{excess noise ratio}:
    \[ F = \frac{\mathrm{ENR}}{Y-1}, \qquad T_e = \frac{T_h - YT_c}{Y-1},
       \qquad F = 1 + \frac{T_e}{T_0} \]
  \item Worked (\code{meas.py}): ENR 15 dB with $Y = 10$ dB gives $F = 3.51$, i.e.
    \mk{5.46 dB}, equivalently $T_e = 729$ K. The floor any such measurement sits on is
    $kT_0 = -174$ dBm/Hz.
\end{itemize}

\penalty0\vspace{2pt}

%=======================================================================
\T{8.10 Surveying the Radiation Zones Round a BTS}

\Q{Working principle of a microwave measuring device to map the radiation zones around a GSM BTS}{\m{8}}
{\color{sub}\footnotesize \tP{1} \yr{\textbf{\texttt{82 Bh}}} \gap$\cdot$\gap paired with ``how microwaves are hazardous'' \m{4} from $\S$7.2; the zone formulas are $\S$7.1}

\begin{itemize}
  \item \textbf{The instrument}: a \textbf{broadband EMF survey meter} (field-strength
    meter), which is a calibrated probe plus a display, usually backed by a
    \textbf{spectrum analyzer} when the individual operators' carriers must be separated.
  \item \textbf{Probe}: an \textbf{isotropic (three-axis) probe}, three mutually
    perpendicular dipoles (E field) or loops (H field) whose outputs are squared and summed,
    so the reading does not depend on how the probe is held. Diode or thermocouple detection,
    high-resistance leads so the cable does not perturb the field.
  \item \textbf{Why both E and H must be measured in the near field}: only in the far field
    is $E/H = 377\ \Omega$ and $S = E^2/377$. Closer in, the ratio is not fixed, so the two
    are measured \textbf{separately} and compared against separate limits ($\S$7.1).
  \item \textbf{Working principle}: the probe dipoles develop a voltage proportional to the
    field; a square-law diode detector gives a dc output proportional to \textbf{power
    density}; high-resistance lines carry it to a meter calibrated in V/m, A/m or
    W/m$^2$.
  \item \textbf{Procedure}: zero the meter with the transmitter off, then walk a
    \textbf{radial grid} out from the tower at a fixed height, log position with GPS, and
    plot contours. Compare each contour with the ICNIRP or FCC limit for that frequency, and
    mark the compliance boundary.
  \item \textbf{Worked} (\code{meas.py}, same formulas as $\S$7.1): a 1.3 m panel at
    1800 MHz has $\lambda = 16.7$ cm, so the reactive near field ends at
    $0.62\sqrt{D^3/\lambda} = $ \mk{2.25 m} and the far field starts at
    $2D^2/\lambda = $ \mk{20.3 m}. With 20 W into a 17 dBi sector, $S$ at 50 m is
    0.032 W/m$^2$, against the ICNIRP public limit of $f/200 = $ 9 W/m$^2$: a margin of
    \mk{about 280 times}.
  \item \textbf{Related ask}: 74 Bh's ``antenna as a DUT, choose a measurement tool''
    \m{8} wants a \textbf{VNA} instead ($\S$8.9): return loss and $S_{11}$ give the match and
    bandwidth, and an anechoic range gives the pattern. It is filed in $\S$7.8.
\end{itemize}

\penalty0\vspace{2pt}

%=======================================================================
\T{8.11 PYQ Mapping}

\lead{Theory questions, covered in this document}
\begin{itemize}
  \item \textbf{Microwave vs low-frequency measurements; major measurement parameters}
    $\rightarrow$ $\S$8.1. \yr{\textbf{71 Bh} \m{3}, \textbf{\texttt{79 Bh}} \m{2},
    \textbf{\texttt{81 Bh}} \m{3}}
  \item \textbf{Choose a proper power measuring tool and explain it} $\rightarrow$ $\S$8.2.
    \yr{72 Ma \m{10} (7.5 mW), 74 Ma \m{8} (airport radar), \textbf{73 Bh} \m{8} (low power)}
  \item \textbf{Bolometry; low power measurement} $\rightarrow$ $\S$8.3.
    \yr{\textbf{70 Bh} \m{6}, \textbf{72 Ash} \m{8}, \textbf{79 Ch} \m{10},
    \textbf{80 Ch} \m{5}}
  \item \textbf{Single-bridge bolometer limitations; double-channel bolometer bridge}
    $\rightarrow$ $\S$8.3. \yr{\texttt{80 Ba} \m{2}, \textbf{\texttt{80 Bh}} \m{8},
    \texttt{81 Ba} \m{4}}
  \item \textbf{Calorimetry; static (dry) calorimeter} $\rightarrow$ $\S$8.5.
    \yr{70 Ma \m{2}, 71 Ma \m{5}, \textbf{71 Bh} \m{7}, \textbf{76 Bh} \m{10},
    \textbf{77 Ch} \m{6}, \textbf{\texttt{81 Bh}} \m{6}}
  \item \textbf{Circulating (flow) calorimetry} $\rightarrow$ $\S$8.5.
    \yr{70 Ma \m{8}, \textbf{75 Bh} \m{5}, \texttt{82 Ba} \m{4}}
  \item \textbf{Calorimeter wattmeter} $\rightarrow$ $\S$8.5. \yr{\texttt{80 Ba} \m{6}}
  \item \textbf{Standing waves and VSWR meter} $\rightarrow$ $\S$8.6.
    \yr{\textbf{70 Bh} \m{4}, \textbf{72 Ash} \m{2}}
  \item \textbf{VSWR of a transmitter with VSWR $>$ 10 (double-minimum method)}
    $\rightarrow$ $\S$8.6. \yr{\textbf{\texttt{79 Bh}} \m{6}}
  \item \textbf{Spectrum analyzer} $\rightarrow$ $\S$8.8.
    \yr{\textbf{\texttt{80 Bh}} \m{4}, \textbf{80 Ch} \m{5}}
  \item \textbf{Network analyzer; VNA} $\rightarrow$ $\S$8.9.
    \yr{73 Ma \m{8}, \textbf{78 Ch} \m{5}}
  \item \textbf{Two-port network model and its parameters} $\rightarrow$ $\S$8.9.
    \yr{\textbf{71 Bh} \m{10}}
  \item \textbf{Measuring device to map radiation zones round a GSM BTS} $\rightarrow$
    $\S$8.10. \yr{\textbf{\texttt{82 Bh}} \m{8}}
\end{itemize}

\lead{Syllabus topics with no PYQ in 24 papers}
\begin{itemize}
  \item \textbf{Thermocouples} ($\S$8.4), \textbf{impedance measurement}, \textbf{measurement
    of unknown loads}, \textbf{measurement of reflection coefficient} and \textbf{frequency
    measurement} ($\S$8.7), and \textbf{noise} ($\S$8.9). All six are written up, briefly,
    because the ``choose a tool'' questions can reach any of them.
\end{itemize}

\lead{Numerical content}
\begin{itemize}
  \item \textbf{This chapter has no numerical component}: no paper in 24 sets a measurement
    calculation. The arithmetic above (duty cycle, the two calorimeter laws, the bridge
    substitution, the double-minimum formula, $\lambda_g$ to frequency) exists to make the
    descriptive answers specific, and every number is re-derived in \code{src/meas.py}.
\end{itemize}

\lead{Corrections and source defects found while building this chapter}
\begin{itemize}
  \item \textbf{Karkee's Chapter-8 PDF export is missing seven slides} that the pptx has:
    impedance measurement, the reflectometer and the low and high VSWR methods (deck slides
    23 to 30). The \textbf{double-minimum method}, which is the whole of 79 Bh's 6 marks,
    appears \textbf{only} in the pptx. Figures are pulled from the pptx for this reason.
  \item Deck slide 30, ``Measurement of frequency'', is an \textbf{empty slide}. $\S$8.7 is
    written from Das ch.\,13 instead \added{}.
  \item Das Fig.\,13.21's worked chart example (\,$S = 2$, $d_{min} = 0.2\lambda_g$\,) prints
    $Z_L = Z_0(1.0 + j0.7)$. Its own equations 13.46 to 13.50 give
    $Z_0(1.55 - j0.69)$, and the closed form $Z_0(1-jSt)/(S-jt)$ agrees with that. The
    printed point does sit on the $S = 2$ circle, but at $70.7^\circ$, which needs
    $d_{min} = 0.348\lambda_g$, not $0.200$. \mk{The answer does not follow from its own
    data}; both routes are asserted in \code{meas.py}.
  \item Das \S13.8 writes insertion loss as $10\log(P_0/P_i)$, which is negative for a lossy
    network. Losses are quoted positive: $10\log(P_i/P_0)$. Same slip as the Chapter-3 deck's
    ($\S$3.4).
  \item Deck slide 19 titles the circulating calorimeter ``\textbf{Circular} Calorimeter'';
    no such instrument. It is circulating, or flow.
  \item The deck sets the high-VSWR threshold at $S > 10$, Das at $S > 20$. Both are quoted
    in $\S$8.6 rather than one being picked.
\end{itemize}
